Adjoin the trivial hyperbolic plane $(\nO, \nInf)$ to $\shW$ and couple it to the dying direction by the coupling construction of \cite[Def.~4.1]{ButlerT4}, whose standing hypothesis $w_1(\shL) = 0$ (the dying line is orientable) is met here, the radical line being trivialised globally by $k = d\Phi$ (clause (5) of Fact \claimref{sheaf.schwarzschild-degeneration-boundary-tame}), taking the normalised transversal $u = \partial_t/m$ (so that $k(u) = 1$, as that construction requires), coupling constants satisfying $\mu\nu < 0$, and bump supported inside the horizon. Then the ambient form $h$ is nondegenerate of constant signature $(2,4)$ at every point of $\Mt$, the boundary $E$ included; the subalgebra generated by $\shW$ is canonically the degeneration sheaf; and the ladder pair
\[
k^* = \frac{\nO}{2\nu}, \qquad k k^* + k^* k = 1, \qquad k^2 = (k^*)^2 = 0
\]
exists exactly over $E$. The ascent $\Cl(1,3) \to \Cl(2,4)$, i.e.\ $\Mat_2(\HH) \to \Mat_4(\HH)$, is thereby performed at the Schwarzschild singularity, localised on $E \cong \R \times S^2$. Moreover no shadow wall forms anywhere in the manifold.
\emph{Provenance:} \cite[Thm.~8.12]{ButlerT4}, verbatim up to the recorded notational alignment, with one clause added here: the parenthetical $w_1(\shL) = 0$ and its discharge by $k = d\Phi$ are not in the source's statement, which invokes \cite[Def.~4.1]{ButlerT4} with that hypothesis standing; the clause records that the hypothesis holds on the Schwarzschild data, by clause (5) of Fact \claimref{sheaf.schwarzschild-degeneration-boundary-tame}, so that the import is applied within its hypotheses.
